\begin{enumerate}
    \item Un noyau radioactif est un noyau instable qui se transforme
          spontanément en un autre noyau en émettant un rayonnement ($\alpha$,
          $\beta$ ou $\gamma$). \hfill\textbf{(1pt)}
    \item 90 protons ($Z=90$) et 147 neutrons ($N=A-Z=237−90=147$).
          \hfill\textbf{(1pt)}
    \item $\ch{^{237}_{90}Th -> ^{A}_{Z}Ra + ^{4}_{2}He}$.
          Lois de conservation de Soddy~: \hfill\textbf{(0,25pt)}
          \[
              \begin{cases}
                  237 = A + 4 &\Rightarrow\quad A = 233 \\
                  90 = Z + 2 &\Rightarrow\quad Z = 88
              \end{cases}
          \]

          Donc~: $\mathbox{\ch{^{237}_{90}Th -> ^{233}_{88}Ra + ^{4}_{2}He}}$
          \hfill\textbf{(0,75pt)}
    \item $t_{1/2}=\frac{\ln 2}{\lambda}
          \quad\Rightarrow\quad
          \lambda=\frac{\ln 2}{t_{1/2}}$ \hfill\textbf{(0,5pt)}

          A.N. \ $\lambda=\frac{\ln 2}{18\times 24\times 3600}
          \quad\Rightarrow\quad
          \mathbox{\lambda \approx \qty{4.5e-7}{\per\second}}$
          \hfill\textbf{(0,5pt)}
    \item Loi de décroissance~:
          $N=N_0 \cdot e^{-\lambda t} \quad\Rightarrow\quad
          m=m_0 \cdot e^{-\lambda t}=
          m_0 \cdot e^{-\frac{\ln 2}{t_{1/2}} t}$.
          \hfill\textbf{(0,5pt)}

          $m=1\times e^{-\frac{\ln 2}{18}\times 36}
          \quad\Rightarrow\quad \mathbox{m=\qty{0.25}{\ug}}$
          \hfill\textbf{(0,5pt)}
    \item $m=m_0\cdot e^{-\lambda t}
          \quad\Rightarrow\quad
          t=-\frac{1}{\lambda}\ln\left(\frac{m}{m_0}\right)
          \quad\text{ou}\quad
          t=\frac{t_{1/2}}{\ln 2}\ln\left(\frac{m_0}{m}\right)$
          \hfill\textbf{(0,5pt)}

          A.N. \
          $t=\frac{18}{\ln 2}\times\ln\left(\frac{1}{0,0156}\right)
          \quad\Rightarrow\quad
          \mathbox{t=\qty{108}{jours}}$
          \hfill\textbf{(0,5pt)}
\end{enumerate}
